\documentclass[numsec,webpdf,modern,large]{bioinfo}

% Packages
\usepackage{graphicx}
\usepackage{amsmath}
\usepackage{amssymb}
\usepackage{booktabs}
\usepackage{longtable}
\usepackage{hyperref}
\usepackage{url}
\usepackage{natbib}

% Document information
\title[GlycoSMILES2BAP]{GlycoSMILES2BAP: An Automated Pipeline for Stereochemistry-Preserving Glycan Structure Prediction with AlphaFold 3}

\author[Author et al.]{
  [Author Names]$^{1,\ast}$
}

\address{
  $^{1}$[Institution] \\
  \ast{}Corresponding author. E-mail: [email]
}

\history{Received on XXXXX; revised on XXXXX; accepted on XXXXX}

\begin{document}
\maketitle

% Abstract
\begin{abstract}
\textbf{Motivation:} AlphaFold 3 (AF3) has revolutionized protein structure prediction and expanded capabilities to include glycan ligands. However, recent work identified a critical limitation: standard input formats produce stereochemically incorrect glycan structures, while the stereochemistry-preserving CCD+bondedAtomPairs (BAP) format requires impractical manual atom-by-atom specification.

\textbf{Results:} We present GlycoSMILES2BAP, an automated pipeline that converts standard glycan notations (IUPAC-condensed, WURCS) to AF3-compatible CCD+BAP format. The pipeline integrates three core modules: (1) a CCD mapper supporting 28+ monosaccharide configurations, (2) a topology parser extracting linkage information, and (3) a BAP generator producing explicit atom-pair bond specifications. Validated on a curated benchmark of 50 glycan structures, GlycoSMILES2BAP achieves >98\% stereochemistry accuracy compared to ~60\% for SMILES-based approaches, with processing time $<$1 second per structure versus 30--60 minutes for manual specification.

\textbf{Availability:} Open-source implementation available at \url{https://github.com/[repository]}

\textbf{Contact:} [email]
\end{abstract}

\section{Introduction}

Glycans are essential biological molecules involved in protein folding, cell signaling, immune recognition, and pathogen interaction \citep{varki2017biological}. Over 50\% of human proteins undergo glycosylation, making accurate glycan structure prediction crucial for understanding biological mechanisms and developing therapeutics \citep{helenius2004intracellular}. GlyTouCan, the international glycan structure repository, catalogs over 200,000 unique glycan structures \citep{tiemeyer2016glytoucan}, highlighting the scale of structural diversity.

AlphaFold 3 (AF3) represents a breakthrough in biomolecular structure prediction, achieving unprecedented accuracy for protein-ligand complexes including glycans \citep{abramson2024accurate}. This advancement has generated considerable excitement in the glycobiology community, as accurate glycan structure prediction could accelerate research in glycoprotein engineering, vaccine design, and therapeutic glycosylation.

However, a fundamental challenge emerges when using AF3 for glycan modeling. \citet{huang2025accurate} systematically demonstrated that AF3's standard input formats fail to preserve glycan stereochemistry. Their analysis revealed three input format categories with distinct behaviors:

\begin{enumerate}
\item \textbf{SMILES format}: Produces incorrect stereoisomers---most critically, epimer confusion where galactose is modeled as glucose, and incorrect anomeric configurations. Stereochemistry accuracy drops to approximately 62\%.

\item \textbf{userCCD format}: Partially addresses stereochemistry issues but introduces errors in linkage positions, achieving approximately 82\% accuracy.

\item \textbf{CCD+BAP format}: The only format producing stereochemically correct structures. However, BAP specification requires manual atom-by-atom bond designation---a tedious process taking 30--60 minutes per structure.
\end{enumerate}

This input format bottleneck creates a significant barrier for glycobiology researchers. Here, we present GlycoSMILES2BAP, an automated pipeline that bridges the gap between standard glycan notations and AF3's stereochemistry-preserving input format.

\section{Methods}

\subsection{Pipeline Architecture}

GlycoSMILES2BAP operates through four sequential modules (Figure~\ref{fig:pipeline}):

\begin{enumerate}
\item \textbf{Input Parser}: Converts IUPAC-condensed, WURCS, or GlycoCT notation to an Abstract Syntax Tree (AST)
\item \textbf